\textbf{Пример 1.} $2^3 = 8 \le 10 < 27 = 3^3$, поэтому $C = 2$. Также $3^2 = 9 \le 10 < 16 = 4^2$, поэтому $S = 3$.

\textbf{Пример 2.} Число $63$ на единицу меньше $64 = 4^3 = 8^2$. Поэтому настоящие кубический и квадратный корни из $63$ чуть меньше $4$ и $8$, а их целые части равны $3$ и $7$: $3^3 = 27 \le 63 < 64 = 4^3$ и $7^2 = 49 \le 63 < 64 = 8^2$.
